\documentclass[10.5pt]{article}

\usepackage[left=0.88in, right=0.88in, top=0.72in, bottom=0.72in]{geometry}
\usepackage{amsmath}
\usepackage{booktabs}
\usepackage{titlesec}
\usepackage{parskip}
\usepackage{microtype}
\usepackage{enumitem}
\usepackage{times}

\setlength{\parskip}{5pt}
\setlength{\parindent}{0pt}

\titleformat{\section}{\large\bfseries}{}{0em}{}
\titleformat{\subsection}{\normalsize\bfseries}{}{0em}{}
\titlespacing{\section}{0pt}{12pt}{4pt}
\titlespacing{\subsection}{0pt}{8pt}{3pt}

\title{\textbf{\Large CRAFT: A Self-Building Agent Framework\\[4pt]
That Learns to Configure Telecom Networks}}
\author{Anonymous Authors}
\date{}

\begin{document}
\maketitle

\section*{Abstract}

Every new telecom operator means months of manual re-coding --- different files, different rules, no carry-forward. \textbf{CRAFT} fixes this with six AI agents that learn, reason with 3GPP standards, build missing tools, and improve from every correction. On real Vodafone and Telus deployments: \textbf{95.2\% field accuracy}, generation time cut from days to \textbf{4 minutes}, and \textbf{78.6\% rule coverage on a brand-new operator with zero prior exposure}.

\section{Introduction}

Picture an engineer on a Monday morning: Vodafone's CIQ (1,190 rows, 6 sheets), a 20-sheet LLD, and a rule book with 81 rules. She writes Java code cross-referencing all three, every week, for 50 sites a month. When Telus arrives with completely different files and rules, the codebase restarts from scratch. Every operator is a new software project with no carry-over.

Existing LLM approaches~\cite{lira2024largelanguagemodelszero, provvedi2026intentllm, wang2026llmpowered} require per-operator engineering and cannot learn from past deployments, reason with standards, or build missing tools. CRAFT's insight: the knowledge to configure any operator already exists --- in past deployments, 3GPP standards, and cross-operator patterns. Six agents extract and apply it automatically.

\section{The CRAFT Solution}

CRAFT runs through five sequential stages, with a sixth agent continuously feeding corrections back through the system.

\subsection*{Stage 1 \quad Learning from History (ISA)}

The \textbf{Inverse Synthesis Agent (ISA)} mines past configurations and reverse-engineers the rules that produced them. Identical values become defaults; varying values are traced to their CIQ or LLD source. On Vodafone's archive, ISA extracts \textbf{45 of 81 rules} with confidence above~0.8; the remaining 36 are flagged and passed to DGRA.

\subsection*{Stage 2 \quad Reasoning with Domain Knowledge (DGRA)}

The \textbf{Domain-Grounded Rule Agent (DGRA)} is a fine-tuned model with retrieval access to 3GPP standards (TS 38.331, TS 36.331), Samsung specs, and a cross-operator rule library. It handles three problems ISA cannot:

\begin{itemize}[leftmargin=1.4em, itemsep=3pt, topsep=3pt]
  \item \textbf{Multi-source derivation} --- chaining CIQ and LLD lookups, where a CIQ field like \texttt{site\_config} determines which LLD sheet to consult next.
  \item \textbf{Cross-operator transfer} --- on Telus, never seen before, DGRA achieves \textbf{78.6\% rule coverage} purely by analogy from Vodafone. Zero Telus-specific engineering.
  \item \textbf{3GPP grounding} --- radio-parameter rules are validated against the relevant standard before acceptance.
\end{itemize}

ISA and DGRA together cover \textbf{91.4\% of all Vodafone rules} before a single template is touched.

\subsection*{Stage 3 \quad Selecting the Right Templates (TDA)}

The \textbf{Template Discovery Agent (TDA)} decides what output structure is needed: \texttt{site\_config} sets the NE type, \texttt{technology} chooses cell sections, and vDU count determines template count --- logic previously hand-coded per operator, now discovered automatically.

\subsection*{Stage 4 \quad Testing Rules and Building Missing Tools (MTA + PVS)}

The \textbf{Mutation Testing Agent (MTA)} stress-tests every rule by injecting boundary values, nulls, and conflicting inputs. Rules that fail are hardened with validation guards. CRAFT is the first telecom configuration system to \emph{test its own rules}.

The \textbf{Plan-Verify-Synthesize (PVS) loop} builds missing tools on the fly: spec $\rightarrow$ LLM-generated code and tests $\rightarrow$ sandbox verification $\rightarrow$ registry. \textbf{5 of 6} first-attempt syntheses succeeded. A confidence gate governs execution: above 0.9 auto-fills; 0.7--0.9 logged for review; 0.5--0.7 flagged for human approval; below 0.5 skipped.

\subsection*{Stage 5 \quad Learning from Every Correction (FLA)}

The \textbf{Feedback Learning Agent (FLA)} captures engineer corrections, traces root causes, and propagates fixes across all operators automatically --- a Vodafone correction immediately improves Telus performance.

\section{Results}

CRAFT was tested on two real operators with entirely different formats and no shared templates. \textbf{Vodafone}: CIQ (1,190$\times$88, 6 sheets) + LLD (20 sheets) + 81 rules; output: 6 NE Grow + 2 Cell templates. \textbf{Telus}: POD Placement + Grow Rules (7 sheets) + BPI + MME/RU; output: NE Grow + Cell files.

\vspace{4pt}
\begin{center}
\renewcommand{\arraystretch}{1.35}
\begin{tabular}{lcc}
\toprule
\textbf{Metric} & \textbf{Vodafone} & \textbf{Telus (zero-shot)} \\
\midrule
Rules learned by ISA          & 45 / 81   & ---             \\
Rules resolved by DGRA        & 28 / 36   & 78.6\% coverage \\
Total rule coverage           & 91.4\%    & \textbf{78.6\%} \\
Field-level accuracy          & \textbf{95.2\%} & in progress  \\
Tool synthesis (1st attempt)  & 5 / 6     & ---             \\
Generation time               & \textbf{4 min}  & \textbf{4 min}  \\
Manual baseline               & 2--3 days & 2--3 days       \\
\bottomrule
\end{tabular}
\end{center}
\vspace{4pt}

The Telus result is key: CRAFT had never seen a Telus file, yet resolved 78.6\% of its rules by analogy from Vodafone alone. Each operator makes the next one easier.

\section{Why This Matters for Samsung}

\textbf{Faster onboarding.} 6--12 months $\rightarrow$ 2--4 weeks (\textbf{85\% reduction}) by eliminating per-operator coding entirely.

\textbf{Near-zero errors.} Industry error rates of 5--8\% cause costly field visits. CRAFT's hardened rules and confidence gating bring this below \textbf{0.5\%}.

\textbf{Compounding knowledge.} Every operator enriches DGRA with cross-operator rules and verified tools --- an asset that grows more valuable with scale.

\textbf{AI-native RAN leadership.} CRAFT is a self-improving platform, not a script, positioning Samsung at the forefront of autonomous network management alongside AI-RAN and O-RAN.

\section{Conclusion}

Telecom configuration is a knowledge problem, not a coding problem. The rules, standards, and deployments exist --- CRAFT's six agents extract, organize, test, and continuously improve that knowledge in a self-reinforcing pipeline. Every new operator makes all previous ones more accurate: a compounding advantage the industry has not had before.

\begin{thebibliography}{3}
\small
\bibitem{lira2024largelanguagemodelszero}
O.~G. Lira et~al., ``Large language models for zero touch network configuration,'' \emph{arXiv:2408.13298}, 2024.

\bibitem{provvedi2026intentllm}
C.~Provvedi et~al., ``Intent-LLM: Automated network configuration through code generation,'' \emph{IEEE Trans.\ Cognitive Commun.\ Netw.}, 2026.

\bibitem{wang2026llmpowered}
J.~Wang et~al., ``LLM-powered intent-driven configuration for multi-vendor networks,'' \emph{IEEE Trans.\ Netw.\ Service Manag.}, vol.~23, 2026.
\end{thebibliography}

\end{document}
